% Part: first-order-logic
% Chapter: tableaux
% Section: provability-quantifiers

\documentclass[../../../include/open-logic-section]{subfiles}

\begin{document}

\olfileid{fol}{tab}{qpr}
\olsection{\usetoken{S}{derivability} మరియు పరిమాణీకరణికాలు}

\begin{explain}
  ఈ విభాగంలో స్థాపించే $\Proves$ గురించిన వాస్తవాలను టాబ్లో నియమాలు
  ఇవ్వడం కూడా సంపూర్ణతా సిద్ధాంతానికి అవసరం.
\end{explain}

\begin{thm}
\ollabel{thm:strong-generalization} $\Gamma$లో గానీ $!A(x)$లో గానీ
కనిపించని స్థిరాంకం $c$ అయి, $\Gamma \Proves !A(c)$ అయితే,
$\Gamma \Proves \lforall[x][!A(x)]$.
\end{thm}

\begin{proof}
$\Gamma \Proves !A(c)$ అనుకుందాం; అంటే $!B_1$, \dots, $!B_n
\in \Gamma$ ఉండి, కింది సమితికి ఒక సంవృత !!{tableau} ఉంది:
\begin{align*}
\{ \sFmla{\False}{!A(c)}, &
\sFmla{\True}{!B_1}, \dots, \sFmla{\True}{!B_n} \}.
\intertext{కింది సమితికి కూడా ఒక సంవృత !!{tableau} ఉందని చూపాలి:}
\{ \sFmla{\False}{\lforall[x][!A(x)]}, &
\sFmla{\True}{!B_1}, \dots, \sFmla{\True}{!B_n} \}.
\end{align*}
ఇచ్చిన సంవృత !!{tableau}లో మొదటి పరికల్పన స్థానంలో
$\sFmla{\False}{\lforall[x][!A(x)]}$ను పెట్టి, పరికల్పనల తర్వాత
$\sFmla{\False}{!A(c)}$ను చేర్చండి.
\begin{center}
\begin{tableau}{not line numbering}
  [\sFmla{\False}{\formula{A}(c)}
    [\sFmla{\True}{\formula{B}_1}
      [\vdots
      [\sFmla{\True}{\formula{B}_n},
        [][][]]]]]
\end{tableau}\qquad
\begin{tableau}{not line numbering}
  [\sFmla{\False}{\lforall[x][\formula{A}(x)]}
    [\sFmla{\True}{\formula{B}_1}
      [\vdots
        [\sFmla{\True}{\formula{B}_n},
          [\sFmla{\False}{\formula{A}(c)}
        [][][]]]]]]
\end{tableau}
\end{center}
ముందు పరికల్పనలుగా లభించిన !!{sentence}s అన్నీ కొత్త
!!{tableau} అగ్రభాగంలో ఇంకా లభిస్తున్నాయి కనుక అది సంవృతంగానే ఉంటుంది.
చేర్చిన వరుస $\TRule{\False}{\lforall}$ను సరిగ్గా వర్తింపజేసిన
ఫలితం: !!{constant}~$c$ అనేది $!B_1$, \dots, $!B_n$లలో గానీ
$\lforall[x][!A(x)]$లో గానీ కనిపించదు; అంటే కొత్త !!{tableau}లో
చేర్చిన వరుస పైన ఎక్కడా కనిపించదు.
\end{proof}

\begin{prop}
\ollabel{prop:provability-quantifiers}
\begin{tagenumerate}{prvEx,prvAll}
\tagitem{prvEx}{$!A(t) \Proves \lexists[x][!A(x)]$.}{}
\tagitem{prvAll}{$\lforall[x][!A(x)] \Proves !A(t)$.}{}
\end{tagenumerate}
\end{prop}

\begin{proof}
\begin{tagenumerate}{prvEx,prvAll}
\tagitem{prvEx}{$\sFmla{\False}{\lexists[x][!A(x)]},
  \sFmla{\True}{!A(t)}$కు సంవృత !!{tableau} ఇది:
  \begin{oltableau}
    [\sFmla{\False}{\lexists[x][\formula{A}(x)]}, just = \TAss
      [\sFmla{\True}{\formula{A}(t)}, just = \TAss
        [\sFmla{\False}{\formula{A}(t)}, just = {\TRule{\False}{\lexists}[1]}, close]]]
    \end{oltableau}}{}
\tagitem{prvAll}{$\sFmla{\False}{\formula{A}(t)},
  \sFmla{\True}{\lforall[x][!A(x)]}$కు సంవృత !!{tableau} ఇది:
  \begin{oltableau}
    [\sFmla{\False}{\formula{A}(t)}, just = \TAss
      [\sFmla{\True}{\lforall[x][\formula{A}(x)]}, just = \TAss
        [\sFmla{\True}{\formula{A}(t)}, just = {\TRule{\True}{\lforall}[2]}, close]]]
    \end{oltableau}
  \sourcecorrection{OLTETAB-012}{ఈ రెండవ ఉదాహరణను పరిచయం చేసే
  ఆంగ్ల మూలపు గణిత భాగం చివర అనవసరమైన కామా ఉంది. అది రెండు
  పరికల్పనల సమితిలో మూడవ ఖాళీ అంశాన్ని సూచించేలా కనిపించకుండా,
  కామాను తొలగించాం; ప్రదర్శిత టాబ్లో, సూత్రాలు మారలేదు.}}{}
\end{tagenumerate}
\end{proof}

\end{document}
